\answer{ex:18:1}{(a)~$0.9949\le w\le1.0049$, es decir, $|1-w|\lesssim0.005$. (b)~$K\ge400$ sinapsis.}
\answer{ex:18:2}{$(\omega_R\tau_w)^2=\sqrt{\alpha(\alpha+2+2\varrho)}/\varrho-1$, $\varrho=\tau_m/\tau_w$; $f_R\approx11.1$~Hz; $|Z(\omega_R)|/|Z(0)|=4.6$.}
\answer{ex:18:3}{(a)~$f=1/\bigl(\tau\ln\frac{\mu}{\mu-\theta}\bigr)$; para $1$~Hz hace falta $\mu/\theta-1\approx3.7\cdot10^{-44}$. (b)~$T=\pi\tau/\sqrt\mu$, $f\approx3.2$~Hz: $f\propto\sqrt\mu$.}
\answer{ex:18:4}{El integrador dispara con $f\lesssim17.7$~Hz (filtro pasa bajos); el resonador, solo en la banda de $5.5$--$22$~Hz (filtro pasa banda).}
\answer{ex:18:5}{(a)~$T_{\min}=\frac{\tau}{w-1}\ln\frac{0.5}{0.3}\approx10.2\ms$. (b)~$B>w-1-0.2=0.3$.}
\answer{ex:18:6}{Constante de tiempo del ajuste $\tau/(\eta\langle r^2\rangle)$, $\eta=\tau\ln10/60\,\text{s}\approx3.8\cdot10^{-3}$; con $I\neq0$ el peso se sesga: hay que restar de $e$ una copia de la orden $I/\tau$.}
\answer{ex:18:7}{Ejercicio abierto. El resonador gana solo un factor de $1.35$ a $11$~Hz y pierde un factor de $17$ a $1$~Hz; la ventaja principal de la reconfiguración está en la selección de banda por umbral (ejercicio~\ref{ex:18:4}).}
